\section{Reproduction}\label{app:repro}

All code is TypeScript under \texttt{src/fms/} of the repository (no runtime
dependencies); the benchmark, tests and paper tables are produced by
\begingroup\footnotesize
\begin{verbatim}
npm ci
npm test        # vitest: Hungarian, A*, PP, CBS, TP, auditor, simulator, properties
npm run bench:fms -- main planner scaling ablation battery    # first round, 420 runs
npm run bench:fms -- stress cascade baselines coupling --jobs 10  # 6,180 runs
npx tsx scripts/paper-data.ts  # tables/*.tex and figures/data/*.csv from the results
npm run check:paper            # every number in the text and tables against the results
\end{verbatim}
\endgroup
The committed results were written at commits \texttt{27080ab} (first round) and \texttt{2d4342a} (second round) of the
repository. Since the second, the code under \texttt{src/fms/} gained
diagnostic counters, options that are off by default, comments, and two
corrections that these runs do not exercise (the stall clock now starts when
a task becomes outstanding, and a token-passing robot that takes a task while
standing on its pickup starts loading at once). Every run is deterministic given its seed (the test suite asserts it
for one configuration). The bench commands above overwrite the committed
files; with \texttt{-{}-out <dir>} they write elsewhere, and
\texttt{npm run check:results -{}- <dir>} compares such a re-run with the
committed records field by field, leaving out the four wall-clock fields
(\texttt{wallMs}, \texttt{plannerMs}, \texttt{allocMs},
\texttt{plannerMsPerTick}). Our repeat of all 6{,}600 runs found no other
difference; its summary is the file
\texttt{paper/experiments/repro/repeat.json}. Regenerating the tables from
a re-run changes their planner-time columns. The last command is a guard.
It regenerates every table and figure data file and compares them byte for
byte; it recomputes every summary file from the per-run records; it
recomputes every numeral written in digits in the abstract, the sections,
this appendix and the figure captions from the committed records and from
constants exported by \texttt{src/fms}, and fails if a digit in that text
has no such check; and it checks which data column every plotted curve
reads. Three kinds of numbers are not recomputed but pinned, which means
that changing one fails the guard until a reviewed snapshot
(\texttt{paper/experiments/guard/pins.json}) is updated: numerals inside the
formulas and algorithms of \cref{sec:method}, numbers written as words
that the script does not list as a claim, and the two facts taken from the
git history (\cref{subsec:setup}). The guard does not check the logic of a
sentence, the references to tables and figures, or the drawing of the
figures beyond their data columns. The bounds stated in
\cref{prop:pp,prop:charge}, \cref{alg:astar} and \cref{eq:cost} are
exercised by unit tests (\texttt{src/fms/\_\_tests\_\_/properties.test.ts}).
Token passing is selected with \texttt{coordinator: 'tp'} or
\texttt{'tpts'} in \texttt{SimConfig}. Our adaptation to the setting of
this paper: the token table stores every robot's complete path with a
permanent rest at its end and no horizon, and space-time A* on it is bounded
by 20{,}000 expansions; a trip is planned as an empty leg to a pickup that
must stay free for the loading dwell, the dwell, and a loaded leg ending in a
rest at the dock; a robot without an eligible task drives to its home bay
(or rests where it is if no path exists) and takes no task on the way; a robot below
$\theta_{\mathrm{lo}}$ takes the nearest free charger slot that is not the
end of another path; and the task-swap step of TPTS only commits a swap if
the displaced robot still has a path to its home bay, which keeps every step
conflict-free.
The operator console shows the
simulator at \texttt{/fms} (\texttt{npm run dev}).

\section{Notation and Parameters}\label{app:params}

\Cref{tab:notation} lists the notation, \cref{tab:params} the default
parameters with the values the sweeps change, and \cref{tab:layouts} the
four layouts.

\begin{table}[h]
\centering
\caption{Notation.}
\label{tab:notation}
\small
\begin{tabular}{ll}
\toprule
$G=(V,E)$ & 4-connected grid graph of free cells\\
$\mathcal{P},\mathcal{D},\mathcal{C},\mathcal{H}$ & pickup points, docks, charger slots, home bays\\
$\mathcal{R}$, $r$ & fleet, robot index; $p^r_t$ position, $s^r_t$ state of charge, $h^r$ home bay\\
$\tau_k=(a_k,p_k,d_k)$ & task: arrival time, pickup, dock; $c_k$ completion time\\
$d(u,v)$ & exact shortest-path distance on $G$\\
$\mathcal{B}$, $w$, $h$ & reservation table, planning window, replanning period\\
$A^{\mathcal{B}}(p,g,t_0)$ & arrival time of windowed space-time A* from $p$ to $g$ under $\mathcal{B}$ (\cref{eq:arrival})\\
$\chi_{\mathcal{B}}$ & congestion count along a free-space shortest path\\
$c(r,\tau)$, $\hat a$ & assignment cost and time-to-pickup estimate (\cref{eq:cost})\\
$w_{\mathrm{b}},w_{\mathrm{c}},w_{\mathrm{h}}$ & battery weight, congestion weight, hold penalty\\
$K$ & candidate tasks priced per robot\\
$\kappa_{\mathrm{d}},\kappa_{\mathrm{p}}$ & in-flight capacity per dock / per pickup\\
$\beta_{\mathrm{m}},\beta_{\mathrm{w}},\gamma$ & charge per move, per wait, charging rate per tick\\
$\theta_{\mathrm{lo}},\theta_{\mathrm{hi}},\rho$ & charging thresholds and feasibility reserve\\
$\delta$ & loading / unloading dwell (ticks)\\
\bottomrule
\end{tabular}
\end{table}

\begin{table}[h]
\centering
\caption{Default parameters (\texttt{DEFAULT\_CONFIG} and \texttt{DEFAULT\_BATTERY} in \texttt{src/fms}).}
\label{tab:params}
\small
\begin{tabular}{lll}
\toprule
\textbf{Parameter} & \textbf{Value} & \textbf{Notes}\\
\midrule
window $w$ / period $h$ & 20 / 5 ticks & as in RHCR \citep{li2021rhcr}\\
candidates $K$ & 6 & nearest pool tasks priced per robot\\
pool size & $\max(2m,8)$ & $m$ eligible robots, oldest tasks first\\
$\kappa_{\mathrm{d}}$ / $\kappa_{\mathrm{p}}$ & 3 / 1 & endpoint reservation capacities\\
$w_{\mathrm{b}}$ / $w_{\mathrm{c}}$ / $w_{\mathrm{h}}$ & 1 / 1 / 10 & \cref{eq:cost}\\
$\beta_{\mathrm{m}}$ / $\beta_{\mathrm{w}}$ / $\gamma$ & 0.001 / 0.0002 / 0.005 & charge units per tick\\
$\theta_{\mathrm{lo}}$ / $\theta_{\mathrm{hi}}$ / $\rho$ & 0.2 / 0.9 / 0.05 & \\
initial charge & $\mathcal{U}[0.5,1.0]$ & $\mathcal{U}[0.2,0.4]$ in the battery stress\\
dwell $\delta$ & 2 ticks & loading and unloading\\
CBS node budget & 2000 & per group, then PP fallback\\
stall limit & 1000 ticks & no task progress $\Rightarrow$ run flagged stalled\\
run limit & 20{,}000 ticks & a run that neither finishes nor stalls stops here\\
tasks per run & 200 (100 on \emph{small}; stress: 100) & pickups, docks sampled uniformly\\
seeds & 1--5 & 1--60 stress, 1--30 cascade, baselines, coupling\\
\bottomrule
\end{tabular}
\end{table}

\begin{table}[h]
\centering
\caption{Benchmark layouts. \emph{console} is generated from the zone / shelf-line / dock counts found in the operator console's wireframe labels (zones A--D with 3, 3, 2, 4 shelf lines; two named shipping docks, each a bay of three delivery cells).}
\label{tab:layouts}
\footnotesize
\begin{tabular}{lrrrrrrr}
\toprule
\textbf{Layout} & \textbf{Size} & \textbf{Free cells} & \textbf{Pickups} & \textbf{Docks} & \textbf{Charger slots (stations)} & \textbf{Home bays} & \textbf{Aisle width}\\
\midrule
small & $16\times10$ & 104 & 4 & 2 & 2 (1) & 8 & open\\
warehouse & $52\times19$ & 715 & 100 & 9 & 8 (4) & 48 & 2\\
console & $36\times25$ & 658 & 84 & 6 & 6 (3) & 32 & 2\\
narrow & $52\times15$ & 496 & 80 & 6 & 8 (4) & 32 & 1\\
\bottomrule
\end{tabular}
\end{table}

\section{Extended Results}\label{app:extended}

\begin{table}[t]
\centering
\caption{\textbf{Seed-paired differences.} Mean per-seed relative difference in \%, with a two-sided 95\% Student-$t$ interval over the 5 seeds (4 degrees of freedom). Runs with the same seed share the task stream and initial charge, so differences are paired. Top: throughput and service time of each cost model relative to Hungarian in every cell of \cref{tab:main}; under $\lambda{=}0.15$ with 16 or 32 robots the makespan is set by the last arrival, so the throughput columns of those rows are uninformative. Bottom: the levers of \cref{subsec:ablation} (throughput relative to the default of the same cost model on the warehouse, 32 robots, batch) and windowed CBS (makespan relative to PP on \emph{small}); $^{\ast}$ marks the differences that survive Holm's procedure at 5\% over these fourteen comparisons.}
\label{tab:paired}
\footnotesize
\setlength{\tabcolsep}{2.5pt}
\begin{tabular}{llrcccccc}
\toprule
 & & & \multicolumn{3}{c}{\textbf{Throughput}} & \multicolumn{3}{c}{\textbf{Service time}}\\
\cmidrule(lr){4-6}\cmidrule(lr){7-9}
\textbf{Layout} & \textbf{Load} & $|\mathcal{R}|$ & Greedy & Proxy & \sname & Greedy & Proxy & \sname\\
\midrule
warehouse & $\lambda{=}0.15$ & 8 & $0.0\pm1.8$ & $-0.5\pm0.7$ & $-0.3\pm0.5$ & $-0.9\pm4.8$ & $+0.5\pm2.6$ & $+0.3\pm1.0$\\
warehouse & $\lambda{=}0.15$ & 16 & $0.0\pm0.8$ & $-0.2\pm0.6$ & $+0.3\pm0.9$ & $+1.0\pm6.6$ & $+4.8\pm6.0$ & $-0.4\pm4.9$\\
warehouse & $\lambda{=}0.15$ & 32 & $-0.3\pm0.9$ & $-0.1\pm0.7$ & $+0.1\pm0.6$ & $+0.2\pm1.0$ & $+0.5\pm0.8$ & $0.0\pm1.2$\\
warehouse & batch & 8 & $+0.2\pm1.6$ & $+0.8\pm1.1$ & $+0.4\pm1.8$ & $+0.3\pm0.9$ & $-0.6\pm1.1$ & $-0.8\pm1.0$\\
warehouse & batch & 16 & $+0.9\pm3.1$ & $+0.4\pm2.4$ & $+1.3\pm4.1$ & $+1.2\pm1.7$ & $+0.3\pm1.7$ & $0.0\pm1.2$\\
warehouse & batch & 32 & $-2.7\pm2.8$ & $+0.3\pm4.7$ & $-0.8\pm1.9$ & $+1.2\pm2.9$ & $-0.2\pm1.5$ & $+1.5\pm1.9$\\
console & $\lambda{=}0.15$ & 8 & $0.0\pm2.9$ & $-1.5\pm2.2$ & $-0.3\pm1.8$ & $+0.4\pm3.5$ & $+2.1\pm3.7$ & $+0.7\pm5.2$\\
console & $\lambda{=}0.15$ & 16 & $0.0\pm0.7$ & $0.0\pm0.9$ & $0.0\pm1.6$ & $-1.0\pm2.2$ & $-0.4\pm2.7$ & $-0.2\pm1.9$\\
console & $\lambda{=}0.15$ & 32 & $-0.3\pm0.4$ & $-0.4\pm0.8$ & $-0.1\pm0.6$ & $+0.5\pm1.7$ & $-0.3\pm1.6$ & $+0.2\pm1.0$\\
console & batch & 8 & $+0.7\pm1.5$ & $+0.2\pm4.3$ & $+2.3\pm4.8$ & $-0.4\pm1.4$ & $+0.1\pm2.2$ & $-0.9\pm1.4$\\
console & batch & 16 & $-0.6\pm5.8$ & $-0.3\pm6.4$ & $-1.2\pm4.5$ & $+1.5\pm1.9$ & $+0.9\pm1.8$ & $+0.4\pm1.1$\\
console & batch & 32 & $-2.0\pm4.5$ & $-1.6\pm2.4$ & $-2.3\pm2.0$ & $+0.3\pm1.7$ & $+0.4\pm3.0$ & $+0.3\pm3.3$\\
\midrule
\multicolumn{3}{l}{\textbf{Lever} (warehouse, 32, batch)} & \multicolumn{6}{c}{\textbf{Throughput} vs.\ default}\\
\multicolumn{3}{l}{window $w{=}10$ (\sname)} & \multicolumn{6}{c}{$+6.9\pm1.6$$^{\ast}$}\\
\multicolumn{3}{l}{window $w{=}30$ (\sname)} & \multicolumn{6}{c}{$-6.9\pm3.9$}\\
\multicolumn{3}{l}{dock capacity 2 (\sname)} & \multicolumn{6}{c}{$-20.7\pm3.3$$^{\ast}$}\\
\multicolumn{3}{l}{dock capacity 2 (Hungarian)} & \multicolumn{6}{c}{$-22.5\pm2.2$$^{\ast}$}\\
\multicolumn{3}{l}{no event epochs (\sname)} & \multicolumn{6}{c}{$-1.2\pm5.2$}\\
\multicolumn{3}{l}{no event epochs (Hungarian)} & \multicolumn{6}{c}{$-5.2\pm3.3$}\\
\multicolumn{3}{l}{replan every tick (\sname)} & \multicolumn{6}{c}{$-0.1\pm4.9$}\\
\multicolumn{3}{l}{no battery term (\sname)} & \multicolumn{6}{c}{$+0.1\pm4.5$}\\
\multicolumn{3}{l}{no battery term (Hungarian)} & \multicolumn{6}{c}{$-1.3\pm4.0$}\\
\multicolumn{3}{l}{candidates $K{=}3$ (\sname)} & \multicolumn{6}{c}{$-0.2\pm3.1$}\\
\multicolumn{3}{l}{candidates $K{=}12$ (\sname)} & \multicolumn{6}{c}{$+0.5\pm4.1$}\\
\multicolumn{3}{l}{\textbf{Planner} (small, batch)} & \multicolumn{6}{c}{\textbf{Makespan} of CBS vs.\ PP}\\
\multicolumn{3}{l}{4 robots} & \multicolumn{6}{c}{$-2.8\pm2.1$}\\
\multicolumn{3}{l}{6 robots} & \multicolumn{6}{c}{$-8.2\pm3.6$$^{\ast}$}\\
\multicolumn{3}{l}{8 robots} & \multicolumn{6}{c}{$-8.2\pm4.3$}\\
\bottomrule
\end{tabular}
\end{table}

\begin{table}[t]
\centering
\caption{\textbf{Fleet-size scaling} on the warehouse layout under batch load (200 tasks, 5 seeds); data of \cref{fig:scaling}. Exp./tick: A* node expansions per simulated tick. Planner ms/tick: median over three repeats on an idle host of the mean over seeds, with [minimum, maximum] (\cref{subsec:setup}).}
\label{tab:scaling}
\footnotesize
\setlength{\tabcolsep}{4pt}
\begin{tabular}{rlccccccc}
\toprule
$|\mathcal{R}|$ & \textbf{Alloc.} & \textbf{Throughput} & \textbf{Service time} & \textbf{Utilisation} & \textbf{Wait actions} & \textbf{A* calls} & \textbf{Exp./tick} & \textbf{Planner ms/tick}\\
\midrule
4 & Hungarian & 4.9\stdv{0.2} & 1909\stdv{64} & 0.78\stdv{0.01} & 18\stdv{10} & 3328\stdv{88} & 18 & 0.02\,{\tiny[0.02,\,0.02]}\\
 & \sname & 4.9\stdv{0.2} & 1914\stdv{68} & 0.78\stdv{0.00} & 25\stdv{11} & 4534\stdv{71} & 25 & 0.02\,{\tiny[0.02,\,0.02]}\\
8 & Hungarian & 9.5\stdv{0.4} & 960\stdv{30} & 0.76\stdv{0.01} & 53\stdv{15} & 3962\stdv{157} & 65 & 0.06\,{\tiny[0.06,\,0.06]}\\
 & \sname & 9.5\stdv{0.3} & 952\stdv{31} & 0.76\stdv{0.01} & 52\stdv{14} & 5118\stdv{87} & 72 & 0.06\,{\tiny[0.06,\,0.06]}\\
16 & Hungarian & 17.5\stdv{0.8} & 491\stdv{20} & 0.73\stdv{0.02} & 118\stdv{20} & 5112\stdv{210} & 221 & 0.17\,{\tiny[0.17,\,0.17]}\\
 & \sname & 17.7\stdv{0.6} & 491\stdv{21} & 0.74\stdv{0.01} & 130\stdv{22} & 6249\stdv{119} & 251 & 0.20\,{\tiny[0.20,\,0.20]}\\
24 & Hungarian & 24.5\stdv{2.2} & 340\stdv{13} & 0.73\stdv{0.05} & 196\stdv{21} & 6398\stdv{240} & 481 & 0.36\,{\tiny[0.36,\,0.37]}\\
 & \sname & 25.1\stdv{2.4} & 338\stdv{13} & 0.74\stdv{0.05} & 188\stdv{26} & 7297\stdv{356} & 521 & 0.39\,{\tiny[0.39,\,0.40]}\\
32 & Hungarian & 30.6\stdv{2.5} & 283\stdv{14} & 0.63\stdv{0.04} & 233\stdv{38} & 8101\stdv{442} & 752 & 0.56\,{\tiny[0.56,\,0.57]}\\
 & \sname & 30.3\stdv{2.2} & 287\stdv{12} & 0.64\stdv{0.04} & 274\stdv{60} & 9282\stdv{359} & 834 & 0.63\,{\tiny[0.63,\,0.64]}\\
40 & Hungarian & 33.0\stdv{2.5} & 259\stdv{10} & 0.50\stdv{0.04} & 287\stdv{35} & 9711\stdv{406} & 930 & 0.69\,{\tiny[0.69,\,0.70]}\\
 & \sname & 32.7\stdv{2.1} & 261\stdv{13} & 0.50\stdv{0.03} & 283\stdv{49} & 12702\stdv{621} & 1080 & 0.83\,{\tiny[0.82,\,0.85]}\\
48 & Hungarian & 33.3\stdv{2.0} & 252\stdv{13} & 0.41\stdv{0.02} & 310\stdv{49} & 12021\stdv{779} & 1021 & 0.76\,{\tiny[0.76,\,0.77]}\\
 & \sname & 33.7\stdv{2.5} & 253\stdv{12} & 0.41\stdv{0.02} & 355\stdv{26} & 16229\stdv{859} & 1275 & 0.98\,{\tiny[0.98,\,0.99]}\\
\bottomrule
\end{tabular}
\end{table}

\begin{table}[t]
\centering
\caption{\textbf{Battery stress}: warehouse, 16 robots, batch load, initial charge drawn from $[0.2,0.4]$ so that every robot must charge during the run (5 seeds).}
\label{tab:battery}
\small
\begin{tabular}{lccccc}
\toprule
\textbf{Variant} & \textbf{Throughput} & \textbf{Service time} & \textbf{Charging sessions} & \textbf{Depletions} & \textbf{Final mean charge}\\
\midrule
Hungarian & 16.3\stdv{0.5} & 626\stdv{20} & 23.2\stdv{2.3} & 0 & 0.37\stdv{0.07}\\
Hungarian, $w_{\mathrm{b}}{=}0$ & 16.6\stdv{0.7} & 618\stdv{19} & 22.6\stdv{2.1} & 0 & 0.38\stdv{0.10}\\
\sname & 16.4\stdv{1.0} & 621\stdv{18} & 23.0\stdv{2.7} & 0 & 0.37\stdv{0.12}\\
\sname, $w_{\mathrm{b}}{=}0$ & 16.6\stdv{0.8} & 624\stdv{17} & 22.6\stdv{2.6} & 0 & 0.34\stdv{0.10}\\
\bottomrule
\end{tabular}
\end{table}

\begin{table}[t]
\centering
\caption{\textbf{What \cref{prop:charge} would need.} $D_\sigma$: largest free-space distance from any cell to any charger slot. $W$, $Q$: largest delay of a trip to a pickup, a dock or a slot, and longest slot wait, measured by the auditor over the normal-battery runs of \cref{tab:stress} (all fleet sizes, planners and loads of the layout). Required: the right-hand sides of \cref{eq:reserve} with these $W$, $Q$ and the default drain ($\beta_{\mathrm{m}}{=}0.001$, $\beta_{\mathrm{w}}{=}0.0002$, $\delta{=}2$); the defaults are $\rho{=}0.05$, $\theta_{\mathrm{lo}}{=}0.2$. With $W{=}Q{=}0$ the default reserve would suffice only on \emph{small}; with the measured $W$ and $Q$ it is short by a factor of 6.6 to 18.6, and not one of the 2{,}160 runs satisfies \cref{eq:reserve} with its own $W$ and $Q$.}
\label{tab:reserve}
\footnotesize
\begin{tabular}{lrrrcccc}
\toprule
\textbf{Layout} & $D_\sigma$ & $W$ & $Q$ & \textbf{Req.\ $\rho$} & \textbf{Req.\ $\rho$ ($W{=}Q{=}0$)} & \textbf{Req.\ $\theta_{\mathrm{lo}}$} & \textbf{Min charge seen}\\
\midrule
small & 17 & 246 & 173 & 0.93 & 0.02 & 0.44 & 0.12\\
warehouse & 67 & 104 & 102 & 0.48 & 0.07 & 0.27 & 0.05\\
console & 56 & 80 & 33 & 0.33 & 0.06 & 0.17 & 0.10\\
narrow & 62 & 156 & 92 & 0.62 & 0.06 & 0.31 & 0.04\\
\bottomrule
\end{tabular}
\end{table}


\Cref{tab:paired} gives the seed-paired comparisons used in
\cref{subsec:main,subsec:ablation}; \cref{tab:scaling} lists the data behind
\cref{fig:scaling}; \cref{tab:battery} the battery stress of
\cref{subsec:main}; \cref{tab:reserve} the delay bounds of
\cref{prop:charge}. The per-run records hold the 38 fields of
\texttt{SimMetrics} (10 that identify the configuration and 28 measured quantities) and 13 auditor fields, with the sweep name, the variant label and
the wall-clock time of the run; the first-round records were written before 4 of these fields were added (three counters of the
baselines and of the CBS backend, and the work-trip delay) and lack them.
The records are in the directory
\texttt{paper/experiments/results/} as one JSON file per sweep, and the
mapping from every number in this paper to a results field is in the file
\texttt{paper/arxiv/NUMBERS.md}.
